% Generated draft: requires mathematical and source review.
\documentclass[11pt]{article}
\usepackage[margin=25mm]{geometry}
\usepackage{fontspec}
\setmainfont{Latin Modern Roman}
\usepackage{amsmath,amssymb,mathtools,unicode-math}
\setmathfont{Latin Modern Math}
\usepackage{enumitem,needspace,xurl,hyperref}
\setlength{\parindent}{0pt}
\setlength{\parskip}{5pt}
\setlength{\emergencystretch}{3em}
\newcommand{\N}{\mathbb{N}}
\newcommand{\Z}{\mathbb{Z}}
\newcommand{\Q}{\mathbb{Q}}
\newcommand{\R}{\mathbb{R}}
\newcommand{\C}{\mathbb{C}}
\newcommand{\F}{\mathbb{F}}
\newcommand{\A}{\mathbb{A}}
\newcommand{\PP}{\mathbb{P}}
\newcommand{\E}{\mathbb{E}}
\newcommand{\eps}{\varepsilon}
\newcommand{\dd}{\,\mathrm{d}}
\begin{document}
\section*{UnsolvedMath ID: 30002024}
\textbf{Source problem code:} catalogue entry
\subsection{Definitions and mathematical statement}
The supplied question concerns the real points of a wonderful variety. The record does not specify all hypotheses on the group or explain the word ``canonical.'' The following gives a conditional interpretation using the usual complex algebraic terminology; these conventions require comparison with the cited report.

Under this interpretation, let \(G\) be a connected complex reductive algebraic group. A wonderful \(G\)-variety is a smooth, complete, irreducible complex algebraic \(G\)-variety \(X\) with an open dense \(G\)-orbit whose complement is a union of finitely many smooth \(G\)-stable prime divisors \(D_1,\ldots,D_r\). These divisors have simple normal crossings and nonempty total intersection, and the \(G\)-orbit closures are precisely their partial intersections, with the empty intersection interpreted as \(X\). These are conventional definitions, not additional hypotheses verified from the report.

Let \(\sigma:G\to G\) be an antiholomorphic involutive group automorphism, interpreted as a real form of \(G\). Let \(\mu:X\to X\) be a real structure, meaning an antiholomorphic involution, satisfying
\[
\mu(g\cdot x)=\sigma(g)\cdot\mu(x)
\qquad(g\in G,\ x\in X).
\]
The source further requires \(\mu\) to be the \emph{canonical} such structure. Its construction, existence hypotheses, and any intended uniqueness assertion are not supplied and remain unresolved qualifications.

Define
\[
X(\R)=\{x\in X:\mu(x)=x\},\qquad
G^\sigma=\{g\in G:\sigma(g)=g\},\qquad
G_0^\sigma=(G^\sigma)^\circ,
\]
where the identity component is taken in the usual real Lie group topology. Equivariance makes \(X(\R)\) invariant under \(G_0^\sigma\). For \(x,y\in X(\R)\), write
\[
x\sim y\quad\Longleftrightarrow\quad
\text{there exists }g\in G_0^\sigma\text{ with }y=g\cdot x.
\]
\textbf{Orbit-count question.} For each wonderful \(G\)-variety with the canonical \(\sigma\)-equivariant real structure intended by the source, determine exactly
\[
N(G,X,\sigma,\mu)
=\#\bigl(G_0^\sigma\backslash X(\R)\bigr)
=\#\bigl(X(\R)/{\sim}\bigr).
\]
The record specifies no particular variety or numerical input, so this is a general determination problem. No nonemptiness assumption on \(X(\R)\) is supplied; the displayed cardinality is zero when \(X(\R)\) is empty. The intended restrictions on \(G\), \(\sigma\), and the canonical structure must be recovered before this conditional formulation can be treated as a precise transcription of the report.
\subsection{Short English statement}
Determine the exact number of orbits of the identity component of the real fixed-point group on the real locus of a wonderful variety carrying the source's canonical equivariant real structure. The precise group hypotheses and meaning of ``canonical'' are missing from the supplied record and require review. The supplied open status is dated curation, not an independently verified conclusion.
\subsection{Sources}
\begin{itemize}
\item Supplied local UnsolvedMath registry record, permanent ID \(30002024\), exact source problem code \texttt{catalogue entry}, title ``Real Orbit Counts on Wonderful Varieties,'' in the collection ``Oberwolfach Reports Open Problems.'' No individual proposer is named; the supplied proposed year is 2012.
\item Supplied primary-source citation: \emph{Mini-Workshop: Generalizations of Symmetric Spaces (2013)}, Oberwolfach Reports, pp.\ 1107--1148, DOI \texttt{10.4171/owr/2012/18}. Supplied source URL: \url{https://doi.org/10.4171/owr/2012/18}. The citation's parenthetical year 2013 differs from the supplied proposed year 2012 and the year appearing in the DOI; this discrepancy is unresolved.
\item Supplied literature source, ``Oberwolfach report'': \url{https://doi.org/10.4171/OWR/2012/18}.
\item Supplied literature source, ``Kac--Moody symmetric spaces'': \url{https://arxiv.org/abs/1702.08426}.
\item Supplied literature source, ``Köhl--Vock'': \url{https://doi.org/10.1515/advgeom-2025-0036}.
\item Supplied literature-check date: \texttt{2026-08-21}. This is supplied dated curation, not independently verified. Its assessment labels the problem open and states that subsequent structural work does not provide a general exact orbit count. The background describes this assessment as literature triage derived from the supplied OWR curation and requests independent verification of time-sensitive status and cited claims.
\item The registry labels the statement \texttt{corrected\_verified} and says that the official report confirms the orbit-count question and supplies definitions omitted by extraction. This is a supplied verification claim; the report's defining passages are not reproduced in the record and have not been independently checked here.
\end{itemize}
\end{document}
